%--------------------------------------------------------------------
\subsection{Density-Matrix Time Propagation}
\label{sec:density-matrix-propagation}
%--------------------------------------------------------------------

The time propagation used in this work builds on the Liouville-von
Neumann density-matrix propagation framework introduced in earlier
work\cite{lvnmd} and on the finite-system \texttt{RMG} RT-TDDFT
implementation developed in our previous
work,\cite{jakowski-rmg-tddft-2025} with the extension that the
propagated matrices are now resolved by crystal momentum and collinear
spin channel.  For each $\mathbf k$ and spin $\sigma$, the propagation
is carried out in a finite subspace
$\mathcal A_{\mathbf k}^{\sigma}$ spanned by ground-state Kohn--Sham
Bloch orbitals,
\begin{equation}
  \Psi_{a\mathbf k}^{\sigma}(\mathbf r)
  =
  e^{\imath\mathbf k\cdot\mathbf r}
  u_{a\mathbf k}^{\sigma}(\mathbf r),
  \qquad
  a\in\mathcal A_{\mathbf k}^{\sigma}.
  \label{eq:tddft_active_bloch_basis}
\end{equation}
The one-particle density matrix in this basis is
\begin{equation}
  P_{ab,\mathbf k}^{\sigma}(t)
  =
  \left\langle
    u_{a\mathbf k}^{\sigma}
    \left|
    \hat\rho^{\sigma}(t)
    \right|
    u_{b\mathbf k}^{\sigma}
  \right\rangle_{\rm cell}.
  \label{eq:density_matrix_definition_kspin}
\end{equation}
The initial condition is the ground-state occupation matrix,
\begin{equation}
  P_{ab,\mathbf k}^{\sigma}(0)
  =
  f_{a\mathbf k}^{\sigma}\delta_{ab},
  \label{eq:initial_density_matrix_kspin}
\end{equation}
where $f_{a\mathbf k}^{\sigma}$ is the ground-state occupation of state
$a$ at $\mathbf k$ in spin channel $\sigma$.

The density matrix obeys the Liouville-von Neumann equation
\begin{equation}
  \imath\hbar
  \frac{d\mathbf P_{\mathbf k}^{\sigma}(t)}{dt}
  =
  \left[
    \mathbf H_{\mathbf k}^{\sigma}(t),
    \mathbf P_{\mathbf k}^{\sigma}(t)
  \right],
  \label{eq:lvn_kspin}
\end{equation}
where
\begin{equation}
  H_{ab,\mathbf k}^{\sigma}(t)
  =
  \left\langle
    u_{a\mathbf k}^{\sigma}
    \left|
    \hat H_{\mathbf k}^{V,\sigma}(t)
    \right|
    u_{b\mathbf k}^{\sigma}
  \right\rangle_{\rm cell}
  \label{eq:hamiltonian_matrix_kspin}
\end{equation}
is the matrix representation of the velocity-gauge Kohn--Sham
Hamiltonian of Eq.~\eqref{eq:velocity_gauge_hamiltonian_k}.  At
nonzero $\mathbf k$, both $\mathbf H_{\mathbf k}^{\sigma}$ and
$\mathbf P_{\mathbf k}^{\sigma}$ are generally complex Hermitian
matrices.

Formally, Eq.~\eqref{eq:lvn_kspin} is solved by a unitary
time-evolution operator,
\begin{equation}
  \mathbf P_{\mathbf k}^{\sigma}(t+\Delta t)
  =
  \mathcal U_{\mathbf k}^{\sigma}(t+\Delta t,t)
  \mathbf P_{\mathbf k}^{\sigma}(t)
  \mathcal U_{\mathbf k}^{\sigma\dagger}(t+\Delta t,t),
  \label{eq:density_unitary_update}
\end{equation}
with
\begin{equation}
  \mathcal U_{\mathbf k}^{\sigma}(t+\Delta t,t)
  =
  \exp\!\left[
    \Omega_{\mathbf k}^{\sigma}(t+\Delta t,t)
  \right].
  \label{eq:magnus_propagator_definition}
\end{equation}
The Magnus generator\cite{Magnus-1954} is written as
\begin{equation}
  \Omega_{\mathbf k}^{\sigma}
  =
  \Omega_{1,\mathbf k}^{\sigma}
  +
  \Omega_{2,\mathbf k}^{\sigma}
  + \cdots,
  \label{eq:magnus_series_kspin}
\end{equation}
where
\begin{equation}
  \Omega_{1,\mathbf k}^{\sigma}
  =
  -
  \frac{\imath}{\hbar}
  \int_t^{t+\Delta t}
  \mathbf H_{\mathbf k}^{\sigma}(\tau)\,d\tau,
  \label{eq:magnus_first_term_kspin}
\end{equation}
and
\begin{equation}
  \Omega_{2,\mathbf k}^{\sigma}
  =
  -
  \frac{1}{2\hbar^2}
  \int_t^{t+\Delta t}d\tau_1
  \int_t^{\tau_1}d\tau_2
  \left[
    \mathbf H_{\mathbf k}^{\sigma}(\tau_1),
    \mathbf H_{\mathbf k}^{\sigma}(\tau_2)
  \right].
  \label{eq:magnus_second_term_kspin}
\end{equation}
In the present method, the Magnus expansion is truncated after its
first term, and the time integral of the Hamiltonian is evaluated using
an endpoint-averaged, midpoint-equivalent approximation,
\begin{equation}
  \Omega_{\mathbf k}^{\sigma}(t+\Delta t,t)
  \simeq
  -
  \frac{\imath\Delta t}{2\hbar}
  \left[
    \mathbf H_{\mathbf k}^{\sigma}(t)
    +
    \mathbf H_{\mathbf k}^{\sigma}(t+\Delta t)
  \right].
  \label{eq:midpoint_magnus_kspin}
\end{equation}
This approximation has the time symmetry associated with midpoint or
trapezoidal integration of the first Magnus term.  The detailed
construction of $\mathbf H_{\mathbf k}^{\sigma}(t+\Delta t)$ in the
self-consistent propagation loop is described in Sec.~\ref{sec:implementation}.

Because $\mathbf H_{\mathbf k}^{\sigma}$ is Hermitian, the generator
$\Omega_{\mathbf k}^{\sigma}$ is anti-Hermitian and
$\mathcal U_{\mathbf k}^{\sigma}$ is unitary.  The action of the
unitary transformation on the density matrix can therefore be written as
a Baker--Campbell--Hausdorff nested-commutator series,
\begin{align}
  \mathbf P_{\mathbf k}^{\sigma}(t+\Delta t)
  &=
  e^{\Omega_{\mathbf k}^{\sigma}}
  \mathbf P_{\mathbf k}^{\sigma}(t)
  e^{-\Omega_{\mathbf k}^{\sigma}}
  \nonumber\\
  &=
  \mathbf P_{\mathbf k}^{\sigma}(t)
  +
  \left[
    \Omega_{\mathbf k}^{\sigma},
    \mathbf P_{\mathbf k}^{\sigma}(t)
  \right]
  +
  \frac{1}{2!}
  \left[
    \Omega_{\mathbf k}^{\sigma},
    \left[
      \Omega_{\mathbf k}^{\sigma},
      \mathbf P_{\mathbf k}^{\sigma}(t)
    \right]
  \right]
  + \cdots .
  \label{eq:bch_density_matrix_kspin}
\end{align}
Equivalently, defining
$\mathbf C_{0,\mathbf k}^{\sigma}=\mathbf P_{\mathbf k}^{\sigma}(t)$
and
\begin{equation}
  \mathbf C_{n,\mathbf k}^{\sigma}
  =
  \frac{1}{n}
  \left[
    \Omega_{\mathbf k}^{\sigma},
    \mathbf C_{n-1,\mathbf k}^{\sigma}
  \right],
  \label{eq:bch_recursive_terms}
\end{equation}
one obtains
\begin{equation}
  \mathbf P_{\mathbf k}^{\sigma}(t+\Delta t)
  =
  \sum_{n=0}^{\infty}
  \mathbf C_{n,\mathbf k}^{\sigma}.
  \label{eq:bch_recursive_sum}
\end{equation}
This form applies equally to real $\Gamma$-point matrices and to the
complex Hermitian matrices that arise for general $\mathbf k$ points.
Exact unitary propagation conserves the trace, Hermiticity, and
eigenvalues of each density matrix; therefore the single-particle
occupations are preserved by the continuous evolution.  In practical
calculations these properties are preserved up to the numerical
truncation and self-consistency errors of the propagation.

The propagated density matrices determine the real-space spin densities.
When the propagated subspace is supplemented by frozen occupied states,
the real-space spin density can be written as
\begin{align}
  \rho^{\sigma}(\mathbf r,t)
  &=
  \sum_{\mathbf k}
  w_{\mathbf k}
  \sum_{a,b\in\mathcal A_{\mathbf k}^{\sigma}}
  u_{a\mathbf k}^{\sigma}(\mathbf r)
  P_{ab,\mathbf k}^{\sigma}(t)
  u_{b\mathbf k}^{\sigma *}(\mathbf r)
  \nonumber\\
  &\quad
  +
  \sum_{\mathbf k}
  w_{\mathbf k}
  \sum_{f\in\mathcal F_{\mathbf k}^{\sigma}}
  f_{f\mathbf k}^{\sigma}
  \left|
    u_{f\mathbf k}^{\sigma}(\mathbf r)
  \right|^2 .
  \label{eq:spin_density_reconstruction}
\end{align}
Here $\mathcal F_{\mathbf k}^{\sigma}$ denotes states that are not
included in the propagated active subspace and remain fixed at their
ground-state occupations.  Equivalently, the same density can be written
as a ground-state density plus the active-space density change,
\begin{equation}
  \rho^{\sigma}(\mathbf r,t)
  =
  \rho_{\rm gs}^{\sigma}(\mathbf r)
  +
  \sum_{\mathbf k}
  w_{\mathbf k}
  \sum_{a,b\in\mathcal A_{\mathbf k}^{\sigma}}
  u_{a\mathbf k}^{\sigma}(\mathbf r)
  \Delta P_{ab,\mathbf k}^{\sigma}(t)
  u_{b\mathbf k}^{\sigma *}(\mathbf r),
  \label{eq:density_reconstruction_delta_p}
\end{equation}
where
$\Delta\mathbf P_{\mathbf k}^{\sigma}(t)
=\mathbf P_{\mathbf k}^{\sigma}(t)-\mathbf P_{\mathbf k}^{\sigma}(0)$.
This form is especially useful for connecting the formal density
reconstruction to an active-space implementation with frozen states.

The propagation is self-consistent because the Hamiltonian is a
functional of the evolving density.  The update cycle can be summarized as
\begin{equation}
  \left\{
    \mathbf P_{\mathbf k}^{\sigma}(t)
  \right\}
  \rightarrow
  \left\{
    \rho^{\sigma}(\mathbf r,t)
  \right\}
  \rightarrow
  \hat V_{\rm H}[\rho](t),
  \hat V_{\rm xc}^{\sigma}[\rho^\uparrow,\rho^\downarrow](t)
  \rightarrow
  \left\{
    \mathbf H_{\mathbf k}^{\sigma}(t)
  \right\}
  \rightarrow
  \left\{
    \mathbf P_{\mathbf k}^{\sigma}(t+\Delta t)
  \right\}.
  \label{eq:self_consistent_density_matrix_cycle}
\end{equation}
The two collinear spin channels are propagated as separate density
matrices, but they are not independent: both contribute to the total
density entering the Hartree potential, and the exchange-correlation
potential in each spin channel depends on $\rho^\uparrow$ and
$\rho^\downarrow$.  The spin structure is discussed in more detail in
Sec.~\ref{sec:collinear-spin-extension}.
